% Transcription — fonds Alexandre Grothendieck, shelfmark 161-6, batch 1 (pages 1-20).
% Established from the facsimile archives/batches/161-6/batch-01.pdf (a mirror of
% https://grothendieck.umontpellier.fr/161-6.pdf), per the skill
% transcribe-grothendieck. The batch file carries no cover sheet: PDF page N is
% archivists' page N, as the pencilled numbers bottom-left confirm.
%
% Pass: Opus 5.5 (claude-opus-5-5), 2026-09-22 — first pass, unchecked against the pages by a human.
% Revised 2026-09-23 (Opus 5.5): p. 18 note on the alpha-chain corrected (the
% chain is exact, alpha(R) = alpha^2(Q) = alpha^3(P) = P; no anomaly); chemise
% p. 1 re-read on the facsimile (259 ppi scan) — the name stays \uncertain,
% the note now gives Laborde as an equally possible reading; p. 17 no longer
% called a "first draft" of pp. 3-4 (nothing on the leaves orders them).
%
% What the batch is.
%   1      The folder's own chemise: on the tab, in pencil, « Graphes,
%          icosaèdre » (whence the inventory's title); at the head, a circled
%          « 6 », a word read « Groupes »? (doubtful), and « (Notes sur formes
%          paradoxiques? prêtées à Labarde?) » — the 2nd and 4th letters of
%          the name are loops his hand does not tell apart as a or o, so
%          « Laborde » reads as well; left doubtful. Cf. folder 162-1, the
%          lending register. No \page; the tab's words head the file and a \note
%          records the rest.
%   2, 11, 20  Blank: absent.
%   3-10   A fair copy in ink on ruled paper, headed « L'icosaèdre » in his
%          hand: the regular icosahedron constructed from a pair (u, v) of
%          orthonormal vectors in a Euclidean 3-space — the equilateral
%          spherical triangle of dihedral angle θ, the choice θ = 2π/5, the
%          caps (« calottes »), the lemma that the six neighbours of x_0 form
%          a cap, the theorem that Γ' (rotations preserving I) acts simply
%          transitively on oriented edges, |Γ'| = 60, |Γ| = 120, and the
%          remark that the automorphism groups of the vertex-edge graph and
%          of the face-edge graph are both Γ'. Page 6 is an unfinished matrix
%          check of the lemma; page 10 closes on the wish to characterise Γ'
%          combinatorially.
%   12, 14 Ink scratch pages on the fifth cyclotomic field: γ = ζ + ζ^{-1},
%          γ² + γ - 1 = 0, the Galois action of (Z/5Z)^*, the fixed ring of
%          ζ ↦ ζ^{-1} computed in the basis 1, T, T², T³, then a sketch over
%          Spec Z[1/n] of μ_n^*, its quotient ϑ_n by inversion and the map
%          λ ↦ λ + λ^{-1} to the affine line.
%   15, 13, foot of 12  A pencil run, crossed through by a long diagonal on
%          15 and 13, on order-preserving maps between ordered sets in a
%          category written (OM): K is shown to be the sum of I and J in (OM),
%          then the statement is generalised to restricted infinite products.
%          Page 13 and the two lines at the foot of 12 are written head to
%          tail with respect to the ink; the reading order is 15 → 13 → foot
%          of 12, and the notes say so. The run begins before page 15 (its
%          hypotheses a), b), c) are not in this batch) — perhaps in batch 2.
%   16-19  Leaves from a spiral notebook, landscape, much of 17-19 written
%          sideways: page 16 is a sheet of doodles (flowers labelled corolle,
%          racine, tige, cœur, pétale; tetrahedra; a cube) with a few lines on
%          Hom_G(G/H, E) = E^H; page 17 is the computation of pp. 3-4 in other
%          letters; page 18 proves G ≃ A_5 via an action on 15 objects
%          (pairs of opposite edges) grouped in 5 triples, over a projected
%          icosahedron with labelled vertices A_i, A'_i; page 19 is scattered
%          scratch (permutations of Z/5Z, normalisers, trial relations for α).
%
% Notation fixed once for the batch: E, S, u^⊥, Σ_u; u ∧ v; x(u,v;α); θ = 2π/5;
% v_i, x_i, x'_i = -x_i, u' = -u; I = I_{u,v}; G = O(E), G' = SO(E); Γ, Γ';
% G_C, G'_C; ℓ, ℓ'; γ^w_α. His ∁ is \complement; « ½ direct » is kept as
% written (semi-direct). On the cyclotomic pages ζ is the root of unity; the
% round letter he writes like a gothic ϑ is rendered \vartheta, uncertain.
% (OM) is kept as two letters in parentheses; what they abbreviate is not on
% these pages.
\documentclass[11pt,a4paper]{article}
% Le préambule commun à toutes les transcriptions du fonds Grothendieck.
%
% Il tient en une idée : **l'incertitude se compose, elle ne se résout pas.**
% Une transcription qui lisse un mot douteux a détruit la seule chose pour
% laquelle on la faisait — un lecteur doit pouvoir savoir, à chaque endroit,
% ce qui était sur le feuillet et ce qui a été ajouté. Les six macros qui
% suivent sont donc l'appareil critique, et non de la décoration.
%
% Les mêmes commandes sont reconnues par `scripts/render.mjs`, qui produit la
% vue de lecture du site. Une macro ajoutée ici doit l'être aussi là-bas,
% faute de quoi le rendu échouera bruyamment — ce qui est voulu : un rendu
% silencieusement faux est pire qu'un rendu absent.

\ProvidesPackage{grothendieck}[2026/08/08 Transcriptions du fonds Grothendieck]

\usepackage{amsmath,amssymb,amsthm}
\usepackage{mathrsfs}
\usepackage{stmaryrd}
% Les diagrammes commutatifs sont partout dans ce fonds : c'est la langue
% même de la géométrie algébrique de Grothendieck, et une transcription qui
% les rendrait en prose ne transcrirait rien.
\usepackage{tikz-cd}
% Français par défaut — c'est la langue des feuillets — mais l'anglais est
% chargé aussi : la traduction et le résumé sont des documents anglais, et la
% typographie française (espaces avant les deux-points) y serait une faute.
% Ces documents basculent par \selectlanguage{english} en tête de corps.
\usepackage[english,french]{babel}
\usepackage{fontspec}
\usepackage{geometry}
\geometry{a4paper,margin=25mm,marginparwidth=28mm}
\usepackage{marginnote}
\usepackage{xcolor}
% ulem, not soul. `soul`'s \st and \ul are built on a character-by-character
% scan that XeLaTeX's font machinery breaks: the document fails with an
% undefined control sequence rather than degrading. ulem does the same two jobs
% and is Unicode-engine safe. `normalem` stops it hijacking \emph, which the
% transcriptions use for Grothendieck's own emphasis.
\usepackage[normalem]{ulem}
\usepackage{hyperref}
\hypersetup{colorlinks=true,linkcolor=black,urlcolor=black,pdfborder={0 0 0}}

% --- Métadonnées du lot -----------------------------------------------------
% Reprises telles quelles par le script de rendu, qui les affiche en tête de
% la vue de lecture. Elles fixent ce que le fichier prétend couvrir : sans
% elles, une transcription flotte sans point d'attache dans un fonds de seize
% mille feuillets.

\newcommand{\folder}[1]{\gdef\grothFolder{#1}}
\newcommand{\batch}[1]{\gdef\grothBatch{#1}}
\newcommand{\leaves}[2]{\gdef\grothFirst{#1}\gdef\grothLast{#2}}
\newcommand{\foldertitle}[1]{\gdef\grothFoldertitle{#1}}
\gdef\grothFolder{?}\gdef\grothBatch{?}\gdef\grothFirst{?}\gdef\grothLast{?}\gdef\grothFoldertitle{}

% --- Le résumé --------------------------------------------------------------
%
% Il ouvre la lecture modernisée, sous le titre « Résumé », et il est composé
% en petit. Ce n'est pas un ornement : c'est le seul passage du document qui
% ne soit pas des mathématiques, et un lecteur doit pouvoir le reconnaître
% avant de l'avoir lu — puis le sauter s'il connaît déjà le sujet, ou s'y
% arrêter s'il ne le connaît pas. Le filet qui le suit marque la reprise du
% corps.

\newenvironment{resume}
  {\small\setlength{\parskip}{0.45em}}
  {\par\medskip\hrule\medskip}

% --- Mots-clés ---------------------------------------------------------------
%
% En anglais, en clôture du résumé de la lecture modernisée : le vocabulaire
% moderne sous lequel chercher ce que les feuillets construisent. Le manifeste
% du site (`npm run manifest`) les extrait pour étiqueter la cote entière —
% cette ligne est la source unique des tags, il n'y a pas de fichier à côté.
\newcommand{\keywords}[1]{\par\smallskip\noindent
  {\footnotesize\textsc{Keywords} --- \textit{#1}}\par}

% --- L'appareil critique ----------------------------------------------------

% Le numéro de feuillet, en marge. C'est l'ancre qui permet de régler un
% désaccord sur une formule en regardant *un* feuillet plutôt qu'une chemise
% de six cents. Le site s'en sert aussi pour tourner les pages du fac-similé
% quand on fait défiler la transcription.
\newcommand{\leaf}[1]{%
  \marginnote{\footnotesize\color{gray!70}\textsf{#1}}%
  \label{leaf:#1}\ignorespaces}

% Illisible : on ne devine pas. Un mot inventé qui se lit comme les autres est
% le pire résultat possible.
\newcommand{\ill}{\textcolor{red!60!black}{[\,\ldots\,]}}

% Lecture douteuse : le mot est donné, et signalé comme incertain.
\newcommand{\uncertain}[1]{\uline{#1}}

% Ajout de l'éditeur — une lettre manquante, un mot sous-entendu.
\newcommand{\add}[1]{\textcolor{blue!50!black}{[#1]}}

% Biffé par Grothendieck lui-même. Ce qu'il a rayé fait partie du document :
% une rature dit souvent où la pensée a changé de direction.
\newcommand{\struck}[1]{\sout{#1}}

% Note du transcripteur, sur l'état matériel du feuillet.
\newcommand{\note}[1]{\par\smallskip{\small\color{gray!60!black}\textit{[#1]}}\par\smallskip}

% Note marginale de Grothendieck, distincte du corps du texte.
\newcommand{\marginal}[1]{\par\smallskip\noindent\hspace{1em}%
  {\small\color{gray!60!black}#1}\par\smallskip}

% --- Titre ------------------------------------------------------------------

\AtBeginDocument{%
  \begin{center}
    {\large\bfseries Fonds Alexandre Grothendieck --- cote n\textsuperscript{o}~\grothFolder}\\[2pt]
    {\itshape\grothFoldertitle}\\[4pt]
    {\small feuillets \grothFirst--\grothLast\ (lot \grothBatch)}\\[2pt]
    {\footnotesize\color{gray!60!black}Transcription établie d'après le fac-similé numérisé
     par l'Université de Montpellier.\\
     Les lectures incertaines sont soulignées, les illisibles notées [\,\ldots\,],
     les ajouts entre crochets.}
  \end{center}
  \vspace{1em}}

\endinput


\folder{161-6}
\batch{1}
\pages{1}{20}
\dating{[à partir de 1973-vers 1977] — le groupe « Dossiers rassemblés par Grothendieck sur différentes thématiques » (161-1 à 162-6) est daté [après 1961-vers 1977]}
\watermark{Édition de démonstration}
\foldertitle{Graphes, icosaèdre [etc.] : notes manuscrites (s.d.)}

\begin{document}

\section*{Graphes, icosaèdre}

\note{Ces mots sont écrits au crayon sur l'onglet de la chemise (page 1). En tête de la chemise, un « 6 » cerclé, un mot lu \uncertain{Groupes}, et, entre parenthèses : « Notes sur formes \uncertain{paradoxiques} prêtées à \uncertain{Labarde} » ; dans le nom, la deuxième et la quatrième lettre sont de petites boucles que la main ne distingue pas entre a et o, et « Laborde » se lit aussi bien (« Lobarde », « Loborde » ne sont pas exclus). La chemise ne porte rien d'autre. Les pages 3 à 10 sont une mise au net à l'encre, titrée de sa main « L'icosaèdre » ; les pages 12 à 19 sont des brouillons, dont une suite au crayon sur les ensembles ordonnés (pp. 15, 13 et pied de la p. 12) et quatre feuillets d'un carnet à spirale.}

\section*{L'icosaèdre}

\page{3}
\marginal{[espace vectoriel muni d'une forme quadratique non dégénérée, qu'on supposera définie positive au besoin]}
$E$ espace vectoriel euclidien de dim 3 $/\mathbb{R}$, produit scalaire $x.y$ ($\|x\| = \sqrt{x^2}$).
$S$ la sphère unité $= \{x \in E \mid \|x\|^2 = 1\}$.

Si $u \in S$, $u^{\perp}$ désigne le plan orthogonal ; donc $u + u^{\perp}$ est le plan tangent à $S$ en $u$, tandis que l'espace vectoriel tangent à $u$ en $S$ est can. isom. à $u^{\perp}$.

L'ensemble des directions tangentes en $u$ \struck{sur} la sphère \add{$S$} s'identifie à l'ensemble $u^{\perp} \cap S = \Sigma_u$ des vecteurs unité orthogonaux à $u$.
\note{le $S$ après « sphère » est ajouté au-dessus de la ligne.}

On suppose $E$ orienté, de sorte que à un couple $u \in S$, $v \in \Sigma_u$ est associé un troisième vecteur $w = u \wedge v$, caractérisé par les conditions $w \in S \cap u^{\perp} \cap v^{\perp}$, $(u,v,w)$ orienté en sens direct.
\marginal{On peut dire que le plan $u^{\perp}$ est orienté : \struck{\uncertain{puisque}} la rotation d'angle $\frac{\pi}{2}$ $v \mapsto u \wedge v$ y est définie.}

Le groupe spécial orthogonal $\mathrm{SO}(E)$ opère de façon simplement transitive sur l'ensemble des directions tangentes à la sphère $S$ en ses différents points, i.e. l'ensemble des couples $(u,v) \in S$ avec $u.v = 0$ i.e. $v \in \Sigma_u$. Le groupe orthogonal $\mathrm{O}(E) \simeq \mathrm{SO}(E) \cdot \mu_2$ ($\mu_2$ = homothéties $\pm 1$) opère de façon doublement transitive.
\note{« $\mathrm{SO}(E)$ » est écrit au-dessus d'un mot biffé, suivi d'un signe noirci, avant $\mu_2$ ; on lit « $(u,v) \in S$ » là où l'on attendrait $S \times S$.}

\struck{\ill{}} \add{Les} points sur $S$ du grand cercle passant par $u$ dans la direction $v$, \struck{sont} i.e. de $w^{\perp} \cap S$, sont de la forme
\[
  x(u,v;\alpha) = u \cos\alpha + v \sin\alpha \qquad \alpha \text{ défini mod } 2\pi\mathbb{Z}.
\]
D'autre part, les \struck{autres} directions tangentes à $S$ en $u$, i.e. les points de $u^{\perp} \cap S$, s'écrivent
\[
  v' = v \cos\theta + w \sin\theta .
\]

\note{en marge gauche, une figure : l'arc de grand cercle de $u$ vers $x(u,v;\alpha)$ et celui de $u$ vers $x(u,v';\alpha)$, l'angle $\theta$ entre eux marqué en $u$ et près des sommets, l'angle $\alpha$ au centre ; en bas, les vecteurs $v$, $v'$, $w$ issus du centre, l'angle $\theta$ entre $v$ et $v'$.}
Fixons $\theta$, et calculons les carrés des côtés du triangle de sommets $u$, $x(u,v;\alpha)$, $x(u,v';\alpha)$. On a
\[
\begin{cases}
  \|x(u,v;\alpha) - u\|^2 \underset{\text{par raison de sym.}}{=} \|x(u,v';\alpha) - u\|^2 \\ \qquad = (\cos\alpha - 1)^2 + \sin^2\alpha = 2(1 - \cos\alpha) = 4 \sin^2 \frac{\alpha}{2} \\[1ex]
  \|x(u,v,\alpha) - x(u,v',\alpha)\|^2 = \sin^2\alpha \, \|v - v'\|^2 \\ \qquad = \sin^2\alpha \; 4 \sin^2 \frac{\theta}{2} = \Bigl(4 \sin^2 \frac{\alpha}{2} \cos^2 \frac{\alpha}{2}\Bigr)\Bigl(4 \sin^2 \frac{\theta}{2}\Bigr)
\end{cases}
\]

\page{4}
Donc pour que ces trois points forment un triangle équilatéral (il est de toutes façons isocèle) il faut qu'on ait
\[
  \text{(1)} \qquad \cos\frac{\alpha}{2} = \pm \frac{1}{2 \sin \theta/2} .
\]
\note{un mot biffé, illisible, après la formule (1).}
\marginal{[\textbf{NB} Un tel $\alpha$ existe sss $|\sin\frac{\theta}{2}| \geqslant \frac{1}{2}$, i.e. si $0 \leqslant \theta \leqslant \pi$, sss $\frac{\pi}{6} \leqslant \frac{\theta}{2} \leqslant \frac{5\pi}{6}$ i.e. \struck{$\frac{\pi}{3} \leqslant$} $\boxed{\theta \geqslant \frac{\pi}{3}}$]}
ou encore, comme $\cos\alpha = \cos^2\frac{\alpha}{2} - \sin^2\frac{\alpha}{2} = 2\cos^2\frac{\alpha}{2} - 1$
\[
  \text{(1')} \qquad \cos\alpha = \frac{1}{2\sin^2\theta/2} - 1 = \frac{1 - 2\sin^2\theta/2}{2\sin^2\theta/2} = \frac{\cos\theta}{2\sin^2\theta/2} = \frac{\cos\theta}{1 - \cos\theta}
\]
\marginal{\textbf{NB} $\begin{cases} \cos 2\alpha = -\frac{3}{5} \\ \cos\alpha = \frac{1}{\sqrt{5}} \end{cases}$ pour $\theta = \frac{2\pi}{5}$}
Bien entendu, dans ce cas, les autres « angles dièdres » définis par ce triangle en les deux autres sommets, sont également égaux à $\theta$. On a en fait construit ici, \add{choisissant la détermination principale de $\alpha$ (comme $0 < \alpha < \pi$),} \emph{le} seul triangle équilatéral inscrit dans la sphère $S$, d'angle dièdre $\theta$, ayant un sommet en $u$, un \struck{autre} \add{deuxième sommet $x$} dans la direction tangente $v$ à $S$ en $u$, et \struck{situé} de façon que \uncertain{pour} le troisième sommet $x'$ l'orientation $\uncertain{u},x,x'$ soit directe.
\note{« choisissant la détermination principale de $\alpha$ (comme $0<\alpha<\pi$) » et « deuxième sommet $x$ » sont des ajouts interlinéaires de sa main.}

Nous allons fixer par la suite
\[
  \text{(2)} \qquad \theta = \frac{2\pi}{5}
\]
et, pour quelque temps, on fixe $u$ et $v$. On posera
\[
  \text{(3)} \qquad v_i = v \cos i\theta + w \sin i\theta \qquad i \in \mathbb{Z}/5\mathbb{Z}
\]
(de sorte que $v_0 = v$, i.e. $v_i = v$ si $i \equiv 0$ (5)), et on posera
\struck{\ill{}}
\[
  \text{(4)} \qquad x_i = x(u, v_i, \alpha) = u\cos\alpha + v_i \sin\alpha = u\cos\alpha + v\cos i\theta \sin\alpha + w \sin i\theta \sin\alpha \qquad (i \in \mathbb{Z}/5\mathbb{Z})
\]
\marginal{\textbf{NB} Il est immédiat que les $u, x_i$ sont \uncertain{6} pts \emph{distincts} (leurs projections sur le plan $u^{\perp}$ sont $0$ et les $v_i \sin\alpha$) à coordonnée suivant $u > 0$ ; donc les $u, x_i, u', x'_i$ sont 12 points \emph{distincts} de $S$}
où $\alpha$ est fixé par (1) ou (1'), $0 < \alpha < \pi$. On trouve donc cinq points en ordre circulaire sur le cercle de \add{centre $u$} rayon (rectiligne) $4\sin^2\frac{\alpha}{2}$ tracé sur la sphère $S$. On pose de plus
\[
  \text{(5)} \qquad x'_i = -x_i , \quad u' = -u \qquad (i \in \mathbb{Z}/5\mathbb{Z})
\]
\marginal{On appelle cet ens. de 12 pts l'\emph{icosaèdre} défini par $u, v$.}
\note{le « rayon (rectiligne) » $4\sin^2\frac{\alpha}{2}$ est le carré de la distance calculée p. 3, non la distance elle-même ; tel quel sur la page.}

Un ensemble de points $(x,y,z)$ de $S$ \add{définissant un triangle équilatéral} d'angle au centre $\alpha$ est appelé un \emph{triangle icosaédral} ; \add{cela signifie aussi qu'il y a une rotation qui l'amène en $(u,x_0,x_1)$} un ensemble de points $(a, b_0, b_1, \ldots, b_4)$ \add{de $S$} est appelé une \emph{calotte d'icosaèdre} s'il existe une transformation \add{orthogonale} $g \in \mathrm{O}(E)$ qui l'amène en l'ensemble précédent $(u, x_0, x_1, \ldots, x_4)$. Il revient encore au même \struck{que} de dire \add{que son cardinal est 6, et que l'on peut ordonner cet ens. de telle façon que} que les triangles $(a, b_i, b_{i+1})$ pour $0 \leqslant i \leqslant 3$ sont icosaédraux (et alors $(a, b_4, b_0)$ est aussi icosaédral*. Le point $a$ est alors uniquement

\page{5}
\note{les deux premières lignes sont la note appelée par l'astérisque de la p. 4, écrite en tête de page.}
* On a encore que les longueurs des dix vecteurs $a - b_i$ ($0 \leqslant i \leqslant 4$), $b_i - b_{i+1}$ ($0 \leqslant i \leqslant 4$, posant $b_5 = b_0$) sont égales à
\[
  \text{(4)} \qquad \ell = 2 \sin\frac{\alpha}{2}
\]
Notons les longueurs correspondantes des vecteurs $a - b'_i = a + b'_i$
\[
  \text{(4')} \qquad \ell' = 2 \cos\frac{\alpha}{2}
\]
\note{le numéro (4) est déjà celui de la formule qui définit les $x_i$, p. 4 ; tel quel.}
\marginal{notons aussi que $\ell^2$ et $\ell'^2$ sont les solutions de l'équation quadratique $\lambda^2 - 4\lambda + 4\sin^2\alpha = 0$}

déterminé, \struck{et l'ens} on l'appelle le \emph{sommet} de la calotte d'icosaèdre, les autres points les \emph{points périphériques}, \struck{et} sur l'ensemble de ceux-ci (grâce à l'orientation donnée de $S$, \struck{donc} il y a alors un ordre circulaire canonique, qui en fait un torseur sous le groupe $\mathbb{Z}/5\mathbb{Z}$. Le sous-groupe $G_C$ de $G = \mathrm{O}(E)$ qui laisse une calotte $C$ invariante, \struck{\ill{}} a dix éléments, \struck{le sous-g} son intersection $G'_C$ avec $G' = \mathrm{SO}(E)$, est le groupe (can. isom. à $\mathbb{Z}/5\mathbb{Z}$) des rotations d'angle multiple de $2\pi/5 = \theta$ autour du \struck{\ill{}} vecteur $a$. \struck{On obtient un} Les cinq points de $G_C \cap \complement G'_C$, d'ordre 2, et formant avec l'élément neutre \struck{sont} les stabilisateurs dans $G_C$ des cinq points périphériques de la calotte ; l'élément $g_i \in G_C \cap \complement G'_C$ qui laisse $b_i$ invariant est la symétrie p.r. au plan engendré par $a, b_i$. La structure du groupe $G_C$ est le produit ½ direct du \struck{\ill{}} sous groupe $G_C$ (can. isom. à $\mathbb{Z}/5\mathbb{Z} = \mu_5$) et du groupe $\mathbb{Z}/2\mathbb{Z}$, opérant sur $G_C$ par symétrie. Enfin, le groupe $G'_C \simeq \mathbb{Z}/5\mathbb{Z} = \mu_5$ opère de façon simplement transitive sur l'ens. des pts périphériques de la calotte (ce qui redonne une définition de l'ordre circulaire can. sur leur ens.)
\note{« ½ direct » : semi-direct. Dans « du sous groupe $G_C$ » et « opérant sur $G_C$ », on n'aperçoit pas le prime que le sens demande ($G'_C$). La phrase « les cinq points de $G_C \cap \complement G'_C$, d'ordre 2, et formant avec l'élément neutre … » reste inachevée telle que les ajouts interlinéaires l'ont laissée.}

Jusqu'à maintenant le choix $\theta = 2\pi/5$ n'a pas joué de rôle particulier : on aurait pu prendre $\theta = 2\pi/m$, et le groupe des $m$ rotations multiples de $2\pi/m$ aurait remplacé celui envisagé. Cela va changer avec le

\textbf{Lemme.} L'ensemble $(x_0; x_1, u, x_{-1}, x'_2, x'_{-2})$ est une calotte d'icosaèdre de sommet \uncertain{$u$}, l'ordre circulaire canonique étant celui écrit.
\note{la notation $(a; b_0, \ldots)$ désigne $x_0$ comme sommet ; la lettre après « sommet » se lit pourtant $u$.}

\emph{Démonstration.} Faisant la symétrie p.r. \struck{au} plan \struck{engendré} engendré par $u, x_0$ (qui échange $x_1$ et $x_{-1}$, $x'_2$ et $x'_{-2}$) on est ramené, à une question d'orientation près, à montrer que les \struck{trois triangles} longueurs des trois vecteurs $x_0 - x'_2$, $x_{-1} - x'_2$, $x'_2 - x'_{-2}$ sont égales à la (4) --- et où on peut se dispenser de la \struck{dernière} vecteur.
\note{au-dessus de « dernière » biffé, un ajout : \uncertain{vérifier pour} \ill{}.}
\struck{\ill{}} Il sera plus astucieux de montrer que si \struck{\ill{}} on considère l'élément $\gamma^{w}_{\alpha}$ de $G' = \mathrm{SO}(E)$ qui est \struck{l'identité sur $w$}, \struck{et qui dans le plan $(u,v)$} est la rotation d'angle $\alpha$ autour de $w$ et l'élément $\gamma^{u}_{\frac{\theta}{2}}$ qui est la \struck{identité sur \ill{} et qui est} rotation d'angle $-\theta/2$ autour de $u$, et si on pose

\page{6}
$\gamma = \gamma^{w}_{\alpha} \gamma^{u}_{-\theta/2}$, alors $\gamma$ transforme la calotte $(u, (x_i)_{i \in \mathbb{Z}/5\mathbb{Z}})$ en l'\struck{\ill{}} \uncertain{ensemble} dit, plus précisément que l'on a les six relations
\[
  \gamma u = x_0 , \quad \gamma x_0 = x'_{2} , \quad \gamma x_1 = x'_{-2} , \quad \gamma x_2 = x_1 , \quad \gamma x_3 = u , \quad \gamma x_4 = x_{-1}
\]
\note{dans $\gamma x_0 = x'_2$ et $\gamma x_2$, l'indice 2 est récrit par-dessus un autre.}

Faisons le calcul :
\begin{align*}
  \gamma u &= \gamma^{w}_{\alpha}(u) = u\cos\alpha + v\sin\alpha \\
  \gamma v &= \gamma^{w}_{\alpha}\Bigl(v\cos\frac{\theta}{2} - w\sin\frac{\theta}{2}\Bigr) = (-u\sin\alpha + v\cos\alpha)\cos\frac{\theta}{2} - w\sin\frac{\theta}{2} \\
  &= -u\sin\alpha\cos\frac{\theta}{2} + v\cos\alpha\cos\frac{\theta}{2} - w\sin\frac{\theta}{2} \\
  \gamma w &= \gamma^{w}_{\alpha}\Bigl(+v\sin\frac{\theta}{2} + w\cos\frac{\theta}{2}\Bigr) = (-u\sin\alpha + v\cos\alpha)\sin\frac{\theta}{2} + w\cos\frac{\theta}{2} \\
  &= -u\sin\alpha\sin\frac{\theta}{2} + v\cos\alpha\sin\frac{\theta}{2} + w\cos\frac{\theta}{2}
\end{align*}
\note{plusieurs signes et plusieurs $\cos$ / $\sin$ de ces lignes sont repassés à l'encre ; on transcrit la forme finale.}
\[
  \gamma = \begin{pmatrix}
    \cos\alpha & -\sin\alpha\cos\frac{\theta}{2} & -\sin\alpha\sin\frac{\theta}{2} \\
    \sin\alpha & \cos\alpha\cos\frac{\theta}{2} & \cos\alpha\sin\frac{\theta}{2} \\
    0 & -\sin\frac{\theta}{2} & \cos\frac{\theta}{2}
  \end{pmatrix}
\]
\begin{align*}
  \gamma u &= u\cos\alpha + v\sin\alpha = x_0 \qquad \text{OK} \\
  \gamma x_0 &= \gamma u \cos\alpha + \gamma v \sin\alpha = u\Bigl(\cos^2\alpha - \sin^2\alpha\cos\frac{\theta}{2}\Bigr) + v\Bigl(\sin\alpha\cos\alpha + \sin\alpha\cos\alpha\cos\frac{\theta}{2}\Bigr) \\
  &\qquad - w\Bigl(\sin\alpha\sin\frac{\theta}{2}\Bigr)
\end{align*}
Or $x'_2 = -x_2 = -u\cos\alpha - v\cos 2\theta\sin\alpha - w\sin 2\theta\sin\alpha$
\note{au-dessus du signe devant $v$, un petit signe de correction, \ill{}.}

à prouver
\[
\begin{cases}
  -\cos\alpha = \cos^2\alpha - \sin^2\alpha\cos\frac{\theta}{2} \\
  -\cos 2\theta \struck{\sin\alpha} = \struck{\sin\alpha}\cos\alpha\Bigl(1 + \cos\frac{\theta}{2}\Bigr) \\
  \struck{\ill{}} \sin 2\theta = \struck{\sin\alpha} \sin\frac{\theta}{2}
\end{cases}
\]
\[
  \gamma x_1 = \gamma u \cos\alpha + \gamma v \sin\alpha\cos\theta + w\sin\alpha\sin\theta
\]
$=$
\note{le calcul s'arrête là ; le reste de la page est blanc.}

\page{7}
\textbf{Cor.} Soit \struck{\ill{}} $g \in G = \mathrm{O}(E)$ transformant la calotte $(u, (x_i))$ en celle du Lemme. Alors $g$ invarie l'icosaèdre $I = I_{u,v}$.
En effet, il est immédiat que l'icosaèdre associé à cette deuxième calotte est encore $I$.

Soit $\Gamma_I = \Gamma_{u,v}$ le sous-groupe de $G = \mathrm{O}(E)$ qui laisse l'icosaèdre $I = I_{u,v}$ invariant, $\Gamma' = \Gamma \cap G'$ le sous-groupe formé des rotations de $\Gamma$. Il est clair que $\Gamma'$ est d'indice 2 dans $\Gamma$ (N.B. $-1 \in \Gamma \cap \complement\Gamma'$) et que $\Gamma$ est donc le produit
\[
  \Gamma \simeq \mu_2 \times \Gamma'
\]

\textbf{Théorème.} $\Gamma'$ opère de façon transitive sur l'ens. des sommets de l'icosaèdre, $I$, et de façon simplement transitive sur l'ens. des couples $(x, d)$ où $x$ est un sommet de l'icosaèdre, $I$ une arête. $\Gamma$ opère de façon simplement transitive sur l'ens. des triples $(x,y,z)$ de points de $I$ formant une face*.
\note{« $I$ une arête » : tel quel ; le $d$ du couple $(x,d)$ est récrit par-dessus une autre lettre.}
\marginal{* On peut aussi \uncertain{dire} sur l'ens. des triples d'arêtes formant \ill{} --- \ill{} c'est évidemment équivalent.}

On a oublié de préciser que par définition les \emph{arêtes} de l'icosaèdre $I$ sont les 30 segments suivants \struck{de longueur $\ell$ qui joignent} :
\[
\begin{cases}
  (u, x_i), & \struck{\ill{}} \; (u', x'_i) \\
  (x_i, x_{i+1}), & (x'_i, x'_{i+1}) \qquad i \in \mathbb{Z}/5\mathbb{Z} \\
  (x_i, x'_{i-2}) & (x'_i, x_{i-2})
\end{cases}
\]
\marginal{Ce sont aussi les \struck{ensembles} couples $(x,y)$ d'éléments distincts de l'icosaèdre, de distance minimum (laquelle est $\ell$) --- il sera plus facile à vérifier, après transitivité de $\Gamma'$ sur l'ens. des sommets.}

De plus, on appelle \emph{faces} de l'icosaèdre \struck{une partie de 3 éléments} des 20 \struck{ensembles} de trois sommets \struck{qui fait partie des ensembles} suivants (ou les triangles qu'ils déterminent)
\[
\begin{array}{lll}
  (u, x_i, x_{i+1}) & (u', x'_i, x'_{i+1}) & i \in \mathbb{Z}/5\mathbb{Z} \\
  (x_i, x'_{i-2}, \underset{\textstyle x_{i+1}}{\underset{\shortparallel}{x_{i-4}}}) & (x'_i, x_{i-2}, \underset{\textstyle x'_{i+1}}{\underset{\shortparallel}{x'_{i-4}}}) &
\end{array}
\]
\note{les trois ajouts interlinéaires de cette phrase (« des 20 », « de trois », « 20 ») et ses ratures laissent une syntaxe flottante ; on la donne telle quelle.}
\marginal{Ce sont aussi les triples $(x,y,z)$ tels que $(x,y)$, $(y,z)$ et $(z,x)$ soient des arêtes. Il y en a exactement \emph{cinq} contenant \struck{issus du sommet} $u$ comme sommet. Ici aussi, ce sera immédiat à vérifier après la transitivité.}

\emph{Démonstration.} On sait déjà que les points $x_i$ \struck{\ill{}} sont congruents sous $\Gamma'$, et par symétrie les points $x'_i = -x_i$ aussi. De plus le lemme nous montre que $u$ est congruent à $x_0$, \struck{$x_0$ est congru} (donc par symétrie $u'$ l'est à $x'_0$), et \struck{par} $x_0$ à $x'_2$ --- donc tous les \uncertain{$x$} sont \emph{congruents}. \struck{Pour montrer} Montrons que $\Gamma'$ est simplement transitif sur l'ens. des couples $(x, d)$ i.e. des arêtes \emph{orientées} $(x,y)$, \struck{\ill{}} i.e. que si $(x',y')$ en est une autre alors $\exists!\, g \in \Gamma'$, $gx' = x$, $gy' = y$. On peut supposer $x = u$, $y = x_0$, et par la transitivité déjà établie sur les sommets, $x' = u$. Donc \struck{\ill{}} alors $y'$ est un des $x_i$, \struck{et l'on sait} il est clair que le stabilisateur

\page{8}
\marginal{Données équivalentes : $\begin{cases} \text{l'icosaèdre + un de ses sommets} \\ \text{icosaèdre + calotte d'icosaèdre contenue dedans} \\ \text{calotte d'icosaèdre} \\ \text{un couple } (u,v) \text{ modulo rotations sur } v \text{ dans } u^{\perp} \text{ d'angle multiple de } \theta \end{cases}$ [\uncertain{point du} fibré principal sur $S$ de groupe \uncertain{$\mathbb{T}$} \uncertain{quotient} \ill{} de celui formé des $(u,v)$ \ill{}]}
\marginal{\struck{\ill{}} Précisons que tout point de l'icosaèdre est centre d'une calotte d'icosaèdre contenue dans $I$, \struck{déterminée} unique, \struck{\ill{}} la calotte de centre $u$ étant formée de $u$ et des \uncertain{sommets} \emph{reliés} à $u$ par une arête.}
\note{la seconde note marginale, écrite en biais dans le coin supérieur gauche, est appelée par un signe cerclé placé après « de $u$ dans ».}

de $u$ dans $\Gamma'$ est aussi le sous-groupe \struck{de} $G'$ \struck{formé des} des $g \in G'$ qui invarient la calotte $(u, (x_i))$, lequel, nous avons vu, est le groupe $\mu_5$ qui opère de façon simplement transitive sur l'ens. des $x_i$, ok.
Prouvons enfin l'assertion de \add{simple} transitivité de $\Gamma$ sur l'ens. des \struck{faces} triples $(x,y,z)$ qui sont une face. Procédant comme dessus, on peut supposer, utilisant ce qui est déjà démontré, que $x = x' = u$, $y = y' = x_0$, $z = x_1$. Mais il n'y a que deux \struck{triangles} faces \struck{\ill{}} admettant $(u, x_0)$ comme arête, leurs autres sommets sont resp. $x_1$ et $x_{-1}$. Si c'est $x_{-1}$, la symétrie p.r. au plan $u, v$ (qui évidemment invarie l'icosaèdre $I$) transforme $x_{-1}$ en $x_1$ en laissant $u, x_0$. Il reste donc : \struck{vérifier} noter que $u, x_0, x_1$ \struck{\ill{}} étant linéairement indépendants une transformation linéaire qui fixe ces trois points est l'identité.

\textbf{Corollaire.} $\Gamma'$ est d'ordre 60, $\Gamma$ d'ordre 120.

Le premier point provient du fait que l'icosaèdre a 30 arêtes, donc $2 \cdot 30 = 60$ arêtes orientées. Comme $\Gamma \simeq \Gamma' \times \mu_2$, on trouve que $\Gamma$ a 120 éléments. On peut aussi noter qu'il y a 20 faces, chacune ayant \struck{trois} six ordres totaux sur l'ens. de ses sommets, ce qui fait $6 \times 20 = 120$ triples ordonnés faisant face.

\textbf{Remarque.} Les sommets et arêtes de l'icosaèdre forment (avec leur relation d'incidence) un \emph{graphe}, et $\Gamma'$ opère par automorphismes de ce graphe. Montrer que \emph{tout} automorphisme \add{g} du graphe est dans $\Gamma'$. On est ramené à \struck{\ill{}} prouver \struck{pour} qu'un automorphisme du graphe qui fixe $u, x_0, x_1$ \uncertain{est} l'identité. Prouvons qu'il est l'identité sur tous les points $x_i$ (qui peuvent être reliés à \struck{\ill{}} un des sommets --- $u$ en l'occurrence --- par une arête) \add{i.e. tous les pts de $I$} : en effet \struck{\ill{}} pour $x_0, x_1$ il n'y a rien à prouver, or $x_2$ ($x_{-1}$) est l'unique point de $I$ \add{distinct de $x_0$ (resp.\ $x_1$)} qui est relié \struck{par} à la fois à $u$ et $x_1$ ($u$ et $x_0$), donc il est fixé, donc aussi $x_3$ (p. ex. comme
\note{« \add{g} » : un petit $g$ ajouté au-dessus de « automorphisme » ; l'ajout « i.e. tous les pts de $I$ » est écrit au-dessus de la parenthèse, avec un signe d'insertion.}

\page{9}
l'unique autre point de $I$, à part $x_{-1}, x_0, x_1, x_2$, qui soit relié à $u$). Donc \struck{l'un} chacun des points \struck{fixes} de la calotte standard \struck{\ill{}} de centre $u$ est fixe. La même chose est vraie pour les points de la calotte de centre $x_0$, donc aussi [appliquant ces deux triangles $(u, x_i, x_{i+1})$] à ceux de la calotte de centre $x_i$, $i \in \mathbb{Z}/5\mathbb{Z}$. Mais cela donne \emph{tous} les points \add{de $I$} sauf $u'$, or celui-ci doit être fixe aussi comme le complémentaire de \uncertain{l'ens.} précédent !
\note{« appliquant ces deux triangles » : lecture douteuse de la tournure ; « de $I$ » est ajouté au-dessus de la ligne.}

\struck{Ainsi, on} On peut aussi regarder l'ens. des \emph{faces} et l'ens. des arêtes de l'icosaèdre $I$ \add{et leur relation d'incidence} (NB une arête est déterminée par l'ens. des deux faces qui la contiennent), ils forment encore un graphe (dont les « sommets » sont les faces de l'icosaèdre ; \struck{on} \uncertain{ce} serait le graphe des sommets-arêtes du dodécaèdre, \emph{dual} à l'icosaèdre). Je dis que le groupe des automorphismes de ce graphe est encore $\Gamma'$ [NB il est \struck{immédiat} évident que $\Gamma$ y opère de façon fidèle]. On est encore ramené à prouver que tout automorphisme de ce graphe qui \struck{\ill{}} invarie une face donnée et ses trois arêtes, disons \struck{\ill{}} celle $(u, x_0, x_1)$, est l'identité. Or il invarie alors chacun des \add{trois} triangles adjacents (comme les uniques triangles adjacents distincts sur les trois arêtes envisagées), \struck{et invarie donc leurs arêtes (?)} puis les \struck{triangles} faces $(u, x_2, x_3)$ et $(u, x_3, x_4)$ comme les seul couple de faces adjacentes mutuellement, et adjacents resp. à $(u, x_1, x_2)$ et $(u, x_0, x_{-1})$ (dont on sait qu'ils sont invariants). Cela donne l'invariance des \struck{faces \ill{}} arêtes $(u, x_{\uncertain{2}})$ et $(u, x_4)$ comme intersection, donc aussi des arêtes restantes $(x_1, x_2)$ et $(x_0, x_{-1})$ des deux \struck{triangles} faces $(u, x_1, x_2)$, $(u, x_0, x_{-1})$. Donc leurs arêtes sont fixes, \struck{donc} et par symétrie on voit que pour chacune des \struck{triangles} faces adjacentes à $(u, x_0, x_1)$, \struck{les \ill{}} ses faces et leurs arêtes sont fixes. À partir de là la conclusion est triviale
\note{« ses faces et leurs arêtes » : tel quel ; le sens demande « les faces adjacentes et leurs arêtes ».}

\page{10}
On voit ainsi que les propriétés des actions de $\Gamma$\struck{\ill{}} sur les ens. de sommets, arêtes et faces de l'icosaèdre (\struck{ensembles} munis de leurs relations d'incidence mutuelles) sont conséquence de la seule structure combinatoire de ces ensembles (soit \struck{le complexe} sommets-arêtes, soit faces-arêtes, soit même le triple système sommets-arêtes-faces …). Il conviendrait cependant de caractériser le sous-groupe $\Gamma'$ de $\Gamma$ en termes combinatoires. On peut dire p. ex. que c'est le groupe qui \struck{respecte} \add{les} ordres circulaires (\add{sur l'ens. des points périphériques des}) \struck{qu'il y a sur les} calottes \struck{sphériques de l'icosa} \add{de} l'icosaèdre (en termes de la structure combinatoire, on peut définir \emph{deux} « syst. cohérents » d'ordres circulaires sur les calottes \struck{\ill{}}, correspondants aux deux orientations possibles de l'espace ambiant).
\note{la phrase sur les ordres circulaires porte deux couches de corrections ; l'état final retenu ici est « le groupe qui [respecte] les ordres circulaires sur l'ens. des points périphériques des calottes de l'icosaèdre ».}

\page{12}
\[
  \zeta^{3} + \zeta^{-1} + \zeta + \zeta^{-3} = \zeta + \zeta^{2} + \zeta^{3} + \zeta^{4} = -1
\]
L'équation de $\gamma$ \uncertain{est} donc
\[
  \boxed{\gamma^{2} + \gamma - 1 = 0}
\]
donc
\[
  \uncertain{\mathfrak{D}}_5 = \mathbb{Z}[\tfrac{1}{5}][t]/(t^{2} + t - 1)
\]
\[
  \boxed{\gamma = (-1 \pm \sqrt{5})/2}
\]
\marginal{\struck{si 2 inv.} Cette formule garde un sens sur tout anneau où 2 \emph{et} 5 sont inversibles}
\note{la lettre de l'équation, $\gamma$, est tracée comme un $\gamma$ bouclé ; celle de l'anneau quotient, un $\mathfrak{D}$ ou un $\vartheta$ gothique, n'est pas sûre.}
\[
  2\cos\frac{2\pi}{5} = (\sqrt{5} - 1)/2 \qquad (2\gamma + 1)^{2} = 5 \qquad \struck{4\gamma^{2} + 4\gamma = 4}
\]
\[
  \cos\frac{2\pi}{5} = \frac{\sqrt{5} - 1}{4} \qquad \cos\frac{4\pi}{5} = \frac{-\sqrt{5} - 1}{4} \qquad c^{2} + c - 1
\]
\[
  (-1 + \sqrt{5})/2 = \bigl((-1 - \sqrt{5})/2\bigr)^{2} - 2 = \frac{1 + 5 + 2\sqrt{5}}{4} - 2 = -\frac{1}{2} + \frac{1}{2}\sqrt{5}
\]
\note{à gauche, un cercle partagé par cinq rayons, deux points du bord marqués d'un trait, un arc double entre deux rayons à droite : la division du cercle en cinq, sans légende.}
\[
  \zeta \longmapsto \zeta^{2} \qquad \zeta \longmapsto \zeta^{3} \qquad \text{si } \zeta + \zeta^{-1} = t
\]
\[
  \struck{\cos' D, D} \qquad \cos'_{\ill{}}(D, \zeta D) = \operatorname{Tr} \zeta^{2} = \zeta^{2} + \zeta^{-2} = t^{2} - 2
\]
\note{dans $\zeta D$, le $\zeta$ est récrit par-dessus une autre lettre.}
\[
  t', t'' \qquad t'' = t'^{2} - 2 = (t''^{2} - 2)^{2} - 2 \qquad t' = t''^{2} - 2
\]
\[
\begin{array}{cccc}
  \zeta & \zeta^{2} & \zeta^{3} & \zeta^{4} \\
  & & \shortparallel & \shortparallel \\
  & & \zeta^{-2} & \zeta^{-1} \\[1ex]
  \zeta + \zeta^{-1} & \zeta^{2} + \zeta^{-2} & \zeta^{3} + \zeta^{-3} & \zeta^{4} + \zeta^{-4} \\
  \shortparallel & \shortparallel & \shortparallel & \shortparallel \\
  \eta = \eta_1 & \eta_2 & \uncertain{\eta_2} & \eta_1
\end{array}
\]
\note{sous $\zeta^{4} + \zeta^{-4}$, une double flèche $\Leftarrow$ pointe vers la gauche.}

\note{au pied de la page, deux lignes au crayon, tête-bêche par rapport au reste, sans rapport apparent avec le calcul qui précède :}
\begin{enumerate}
  \item[d)] $\gamma((V_\lambda)) \leq \alpha_\lambda(V'_\lambda) \Longrightarrow V_\lambda \leq V'_\lambda$ ou $\exists\, \mu \neq \lambda$ avec $V_\mu = 0_{\uncertain{I_\mu}}$
  \item[e)] $\gamma((V_\lambda)) \subset \sup_i W_i$
\end{enumerate}
\note{le $V$ est tracé comme un $\forall$ barré ; la lecture $V$ est une conjecture.}

\page{13}
\note{page au crayon, écrite tête-bêche par rapport aux pages d'encre qui l'entourent, et barrée d'un long trait oblique ; elle ne traite pas de l'icosaèdre, mais d'applications croissantes entre ensembles ordonnés. Elle commence au milieu d'un calcul dont le début n'est pas dans ce lot, et se termine par « TSVP » : la suite est, selon toute apparence, les deux lignes d) et e) au pied de la p. 12. Les lettres $U$, $V$, $W$ se confondent souvent sous ce crayon.}
\[
\begin{aligned}
  &= \sup_{\ill{}} [\bar\alpha(U) \wedge \bar\beta(V)] \wedge [\bar\alpha(U') \wedge \bar\beta(V')] \\
  &= \bigl(\sup \bar\alpha(U) \wedge \bar\beta(V)\bigr) \wedge \bigl(\sup \bar\alpha(U') \wedge \bar\beta(V')\bigr) \\
  &= \bar\gamma(W) \wedge \bar\gamma(W') , \qquad \text{ok}
\end{aligned}
\]

$\gamma$) $\bar\gamma$ commute aux Sup quelconques
\[
  \bar\gamma(\underbrace{\struck{\ill{}}\sup W_i}_{W}) \geq \sup\bigl(\bar\gamma(W_i)\bigr) \quad \text{clair, prouvons } \leq
\]
\[
  \bar\gamma(W) = \sup_{\gamma(U \times V) \subset \sup_i W_i} \bar\alpha(U) \wedge \bar\beta(V)
\]
\note{la ligne précédente est reconstituée d'un « $\shortparallel$ » placé sous $\bar\gamma(\sup W_i)$, d'un « Sup » et, dessous, de la condition « $\gamma(U \times V) \subset \sup W_i$ ».}
à prouver que l'on a $\bar\alpha(U) \wedge \bar\beta(V) \leq \sup_i \bar\gamma(W_i)$. Or on a, si $\gamma(U \times V) \subset \sup W_i$ :
\[
  \begin{array}{l} U = \sup U_j \\ V = \sup V_k \end{array} \;\Big|\; \forall\, j, k \ \exists\, i \text{ avec } \gamma(U_j \times V_k) \leq W_i
\]
Donc
\[
  \bar\alpha(U) \wedge \bar\beta(V) = \bigl(\sup_j \bar\alpha(U_j)\bigr) \wedge \bigl(\sup_k \bar\beta(V_k)\bigr) = \sup_{j,k} \bar\alpha(U_j) \wedge \bar\beta(V_k) \qquad \Bigl(\overset{?}{\leq} \sup_i \bar\gamma(W_i)\Bigr)
\]
et comme $\exists\, i$ $\surd$ $\gamma(U_j \times V_k) \leq W_i$, on a
\[
  \bar\alpha(U_j) \wedge \bar\beta(V_k) \leq \bar\gamma(W_i) \leq \sup_{i'} \bar\gamma(W_{i'}) \qquad \text{qfd}
\]

Généralisation aux sommes (produits) infinies :
\[
  (I_\lambda)_{\lambda \in \Lambda} \qquad I_\lambda \in \operatorname{Ob}(\uncertain{\mathrm{OM}})
\]
$\prod' I_\lambda$ partie du produit formée des $(V_\lambda)$ tels que $V_\lambda = 1_{I_\lambda}$ pour presque tout $\lambda$
\[
  \alpha^{0}_{\mu} : I_\mu \longrightarrow \prod{}' I_\lambda \qquad \alpha^{0}_{\mu}(x)_\lambda = \begin{cases} x & \text{si } \lambda = \mu \\ 1_{I_\lambda} & \text{si } \lambda \neq \mu \end{cases} \qquad \text{appl. cr.}
\]
\[
  \gamma : \prod{}' I_\lambda \longrightarrow K \quad \text{flèche dans } (\uncertain{\mathrm{OM}})
\]
\[
  \alpha_\mu = \gamma \circ \alpha^{0}_{\mu} : I_\mu \to K \quad \text{appl. cr.}
\]
On suppose
\begin{enumerate}
  \item[a)] Les $\alpha_\mu$ commutent aux Sup quelc.
  \item[b)] $\gamma$ commute aux Inf finis [en particulier $\gamma((V_\lambda)) = \inf_\lambda \alpha_\lambda(V_\lambda)$ (Inf en fait fini) ; et les $\alpha_\mu$ sont des morphismes de $(\uncertain{\mathrm{OM}})$]
  \item[c)] Tout élément de $K$ est le sup d'éléments de $\prod'_\lambda I_\lambda$
\end{enumerate}
TSVP
\note{« (OM) » : deux lettres capitales entre parenthèses, désignant une catégorie d'ensembles ordonnés ; la lecture des lettres n'est pas sûre. Au c), on attendrait « d'éléments de $\gamma(\prod' I_\lambda)$ ».}

\page{14}
\note{page de calculs épars, sans texte suivi ; on les donne dans l'ordre de lecture, de haut en bas et de gauche à droite. La lettre ronde tracée comme un $\vartheta$ (déjà p. 12) est rendue $\vartheta$ ; le $\mu$ à double jambage est rendu $\mu$.}
\[
  \Gamma = (\mathbb{Z}/5\mathbb{Z})^{*} \simeq \mathbb{Z}/4\mathbb{Z}
  \qquad
  \begin{array}{ccc}
    \mu_5 & \supset & \mu_5^{*} \\
    | & & | \\
    \operatorname{Spec} \mathbb{Z} & \longleftarrow & \mathbb{Z}[\frac{1}{5}]
  \end{array}
\]
fini étale de rang 4 sur $\mathbb{Z}[\frac{1}{5}]$ ; \struck{\ill{}} Galoisien de groupe $\Gamma$.
\[
  \struck{\ill{}}\; \tau = \mu_5^{*}/\beta \quad \text{engendré par } \struck{\zeta}\, \zeta + \bar\zeta = \zeta + \zeta^{-1} = \eta
\]
\note{sous le premier $\zeta$, une flèche : « racine primitive universelle ».}
\begin{itemize}
  \item deux générateurs $(\alpha, \alpha')$ : 2 et 3 ($= -2$) mod 5
  \item un élément d'ordre 2 ($-\beta$) : $-1$ ou 4 mod 5
  \item un élément d'ordre 1 : 1 mod 5
\end{itemize}
Sur $\mu_5(S)$ l'effet est
\[
  \zeta \mapsto \zeta^{2} \text{ ou } \zeta \mapsto \zeta^{-2} = \zeta^{3} \text{ resp.} \qquad \zeta \mapsto \zeta^{-1} = \zeta^{4} \qquad \zeta \mapsto \zeta
\]
\[
  \begin{array}{c} B \supset B^{\Gamma} \\ | \\ A \end{array}
  \qquad
  B = A[\zeta] \struck{\ill{}} = A[T]/(1 + T + T^{2} + T^{3} + T^{4})
  \qquad T \longmapsto T^{4}
\]
\[
\begin{aligned}
  \beta(T) &= T^{-1} = T^{4} = -(1 + T + T^{2} + T^{3}) \\
  \beta(T^{2}) &= T^{-2} = T^{3} \\
  \beta(T^{3}) &= T^{-3} = T^{2} \\
  \beta(T^{4}) &-
\end{aligned}
\qquad\qquad 1 = T^{0} \mid T \quad T^{2} \quad T^{3}
\]
\[
\begin{aligned}
  \beta(c_0 + c_1 T + c_2 T^{2} + c_3 T^{3}) \struck{\ill{}}
  &= c_0 \mathbin{\ill{}} c_1(\uncertain{-1} + T + T^{2} + T^{3}) + c_2 T^{3} + c_3 T^{2} \\
  &= (c_0 - c_1) - c_1 T + (c_3 - c_1) T^{2} + (c_2 - c_1) T^{3} \\
  &\overset{?}{=} c_0 + c_1 T + c_2 T^{2} + c_3 T^{3}
\end{aligned}
\]
\[
  c_1 = 0 \qquad c_2 = c_3 \qquad\qquad c_0 + c_2 (T^{2} + T^{3})
\]
\note{le signe entre $c_0$ et $c_1$ est couvert d'une tache ; la ligne suivante suppose $c_0 - c_1(1 + T + T^{2} + T^{3})$.}

${}_n\mathcal{U}$ étale de \uncertain{rang} \ill{} sur $\mathbb{Z}$ si $n$ \uncertain{premier} ; \struck{\ill{}} $\mu_n$ \emph{tordu} par $T$
\[
  \begin{array}{c} \mu_n \\ \cup \\ \mu_n^{*} \\ \downarrow \\ S \end{array}
  \quad (\mathbb{Z}/n\mathbb{Z})^{*} \qquad \zeta \longmapsto \zeta^{-1}
\]
\note{à gauche, un grand cercle sans légende, marqué $\varphi(n)$ sur son bord droit ; deux traits obliques relient $\mu_n^{*}$ à $S$.}
\[
  \varphi(n) = \prod \bigl(p_i^{\alpha_i} - p_i^{\alpha_i - 1}\bigr) \qquad \text{pas gros} \qquad p^{n} - p^{n-1} = p^{n-1}(p - 1) \qquad (\mathbb{Z}/n\mathbb{Z})^{*} \ni -1
\]
\uncertain{or}
\[
\begin{array}{ccccc}
  \mu_n^{T} & \subset & \mathcal{U} & & \\
  \cup & & & & \\
  \mu_n^{*T} & \subset & \mathcal{U} & & \zeta + \bar\zeta \\
  \Big\downarrow{\scriptstyle \text{étale rang } 2} & & \Big\downarrow{\scriptstyle \operatorname{Tr}_{K/\mathbb{Q}}} & & \\
  \vartheta_n^{T} = \mathcal{V}_n & \subset & \mathcal{O} & & \\
  \Big\downarrow{\scriptstyle \text{étale rang } \uncertain{\varphi(n)/2}} & & & & \\
  S & & & &
\end{array}
\]
\[
  \begin{array}{ccc}
    \mu_n & \subset & \mathbb{G}_m \\[1ex]
    \mu_n^{*} & \subset & \mathbb{G}_m \\
    \downarrow & & \downarrow{\scriptstyle \rho} \\
    \vartheta_n & \subset & \mathbb{E}^{1}
  \end{array}
  \qquad \rho(\lambda) = \lambda + \lambda^{-1}
\]
\[
  \zeta + \zeta^{-1} \;\Big|\; \begin{array}{l} \frac{p - 1}{2} = 2 \\ \cos \end{array}
\]
\[
  1 + \zeta + \zeta^{2} + \zeta^{3} + \zeta^{4} = 0 \qquad \zeta + \zeta^{-1} = \gamma
\]
\[
\begin{aligned}
  \operatorname{Tr} \gamma &= \gamma + \gamma^{\alpha} = \zeta + \zeta^{-1} + \zeta^{2} + \zeta^{-2} = \zeta + \zeta^{2} + \zeta^{3} + \zeta^{4} = -1 \\
  \operatorname{N} \gamma &= \gamma\gamma^{\alpha} = (\zeta + \zeta^{-1})(\zeta^{2} + \zeta^{-2}) =
\end{aligned}
\]
\note{le calcul de la norme s'arrête sur ce signe $=$.}

\page{15}
\note{au crayon, barrée d'un long trait oblique comme la p. 13. Cette page précède en fait celle-là : son dernier calcul, $\bar\gamma(W \wedge W') = \ldots$, se poursuit en tête de la p. 13, qui continue par le point $\gamma$) après les $\alpha$) et $\beta$) d'ici. Le début (hypothèses a), b), c) et la catégorie (OM)) n'est pas dans ce lot.}
Alors $\alpha, \beta : \begin{array}{l} I \to K \\ J \to K \end{array}$ [font de $K$ une somme de $I, J$ dans $(\mathrm{OM})$]
\marginal{sont dans $(\mathrm{OM})$ (grâce à a) b)), je dis qu'ils}
\note{la ligne interlinéaire au-dessus, reliée par un trait, s'insère après « $J \to K$ ».}

a) $\alpha, \beta$ épimorphiques i.e. pour $L \in \operatorname{Ob}(\mathrm{OM})$,
\[
  \struck{\ill{}} \qquad \operatorname{Hom}_{\mathrm{OM}}(K, L) \longrightarrow \operatorname{Hom}_{\mathrm{OM}}(I, L) \times \operatorname{Hom}_{\mathrm{OM}}(J, L) \qquad \varphi \longmapsto (\varphi \circ \alpha, \varphi \circ \beta)
\]
injectif.

En effet, si on connaît $\varphi \circ \alpha$ et $\varphi \circ \beta$, i.e. $\varphi$ sur les $\gamma(x \times 1)$ et $\gamma(1 \times y)$, on les connaît sur leurs $\wedge$ deux à deux $\gamma(x \times y)$ (par b)) donc sur les Sup de \uncertain{objets} de \uncertain{ces} $\gamma(\ldots)$ donc sur tout $K$ (par c)).

b) \uncertain{Ceci} est surjectif. Donnons-nous donc
\[
  \bar\alpha : I \longrightarrow L \qquad \bar\beta : J \longrightarrow L \qquad \text{dans } (\mathrm{OM})
\]
définissons une application
\[
  \bar\gamma : K \longrightarrow L
\]
par
\[
  \bar\gamma(W) = \sup_{\substack{U \in I \\ V \in J \\ \gamma(U \times V) \subset W}} \bar\alpha(U) \wedge \bar\beta(V)
\]

1°) Prouvons d'abord
\[
  \bar\gamma \circ \alpha = \bar\alpha , \qquad \bar\gamma \circ \beta = \bar\beta
\]
Or
\[
  (\bar\gamma \circ \alpha)(U') = \bar\gamma(U' \times 1) = \sup_{\gamma(U \times V) \leq \gamma(U' \times 1)} \bar\alpha(U) \wedge \bar\beta(V)
\]
Or les $U \times V$ en question sont, en vertu de d), ceux pour lesquels ou bien $V = 0_I$ (l'élément correspondant du Sup est alors $0_{\uncertain{K}}$, et \uncertain{ne} compte pas), ou bien ceux pour lesquels $U \subset U'$, $V$ quelconque. Donc le Sup est $\bar\alpha(U') \wedge \bar\beta(1) = \bar\alpha(U')$, OK. On prouve de même $\bar\gamma \circ \beta = \bar\beta$.
\note{« en vertu de d) » : la condition d) est l'une des deux lignes au crayon du pied de la p. 12. $0_I$ tel quel (on attendrait $0_J$).}

2°) Prouvons que $\bar\gamma$ est un morphisme de $(\mathrm{OM})$.
\begin{enumerate}
  \item[$\alpha$)] $\bar\gamma$ croissant : clair
  \item[$\beta$)] $\bar\gamma$ commute aux Inf finis.
\end{enumerate}
\[
  \bar\gamma(1_K) \struck{\ill{}} \overset{\text{def}}{=} \bar\alpha(1_I) \wedge \bar\beta(1_J) = 1_K
\]
\marginal{Inf de famille vide}
\[
  \bar\gamma(W \wedge W') = \sup_{\gamma(U \times V) \subset W \wedge W'} \bar\alpha(U) \wedge \bar\beta(V) = \sup_{\substack{\gamma(U \times V) \subset W \\ \gamma(U' \times V') \subset W'}} \bar\alpha(U \wedge U') \wedge \bar\beta(V \wedge V')
\]
[d'où $\gamma(U \wedge U' \times V \wedge V') \subset W \wedge W'$ via (b)]
\note{« $= 1_K$ » : tel quel ; on attendrait $1_L$. Les lettres $U$, $V$ sont peu distinctes au crayon.}

\page{16}
\note{feuillet de carnet à spirale, en largeur, couvert de croquis à l'encre et au crayon : fleurs à trois, quatre ou cinq pétales, dont l'une est légendée, en marge gauche, « corolle », « racine », « tige », « cœur », « pétale » ; tétraèdres et pyramides, dont l'un avec un point intérieur relié aux sommets ; un losange à diagonales suivi d'un « $G$ » ; un hexagone découpé en six triangles, relié par un arc au sommet d'une fleur et légendé « $\mathrm{O}(3)$ », « $\mathrm{SO}(3)$ » ; une suite d'ovales emboîtés dont les deux plus intérieurs portent « $u_1$ » et « $u_0$ » cerclé ; un carré inscrit dans un carré, reliés par les diagonales ; un graphe sur six sommets marqués de ronds et de croix ; un cube dont trois faces sont fléchées vers un sommet. On transcrit ci-dessous les formules et les mots, sans restituer la disposition.}
\[
  \underline{X} \longrightarrow (\mathrm{Ens}) \qquad \underline{\hat X} \longrightarrow (\mathrm{Ens}) ,
\]
\uncertain{correspondant} \ill{}
\[
  x \subset e \qquad x \quad y \qquad z \to x \cap y
\]
\[
  4 \longrightarrow 3 \longrightarrow 2 \longrightarrow 1 \longrightarrow 0
\]
\[
  X_1 \longleftarrow X_0 \qquad y \leq x \qquad U_x \subset U_y \qquad F(y) \longrightarrow F(x)
\]
\note{en haut, au crayon, « $8 \times 6$ » et, à côté d'un symbole inachevé, « \uncertain{789} » ou « \uncertain{989} ».}

\uncertain{opinion} \ill{} \ill{} \uncertain{homogène} $X \simeq G/H$ \ill{} $= [x \in S^{2}]$ et \uncertain{même} cas \uncertain{\ill{}} \uncertain{particulier} de $S^{1}$).
\[
  \operatorname{Hom}_G(X, E) \subset \operatorname{Hom}_G(P_S, E) \simeq E
\]
\uncertain{savoir} l'unique application $g \longmapsto gx$ telles que l'on ait $gx = ghx$ i.e. $E^{H}$
\note{« $P_S$ » : lecture douteuse. Le passage « $gx = ghx$ i.e. $E^{H}$ » décrit, semble-t-il, l'identification de $\operatorname{Hom}_G(G/H, E)$ aux invariants $E^{H}$.}

\page{17}
\note{feuillet de carnet à spirale, écrit en travers, au crayon et à l'encre : le calcul des pp. 3--4, avec d'autres lettres ($\rho_i$, $v'$, $w'$, $x(\alpha)$, $y(\alpha)$). Des croquis l'accompagnent : un cône de sommet le centre, découpé en triangles d'angle au sommet $\alpha$ (côtés marqués $1$, $\rho$, $\alpha$) et légendé « $0 \leqslant i \leqslant 4$ » ; les vecteurs $u$, $v$ issus du centre d'une sphère ; deux cercles tangents, l'un portant deux axes, l'autre un angle.}
$u, v, w$
\[
  \struck{\ill{}} + v \cos i\frac{2\pi}{5} + w \sin i\frac{2\pi}{5} = \rho_i \qquad \boxed{0 \leqslant i \leqslant 4}
\]
\[
\begin{cases}
  u \\
  u\cos\alpha + \rho_i \sin\alpha \\
  -u\cos\alpha - \rho_i \sin\alpha \\
  -u
\end{cases}
\]
\note{le $u$ de la première ligne de l'accolade est écrit par-dessus un gribouillis qui couvre aussi le début de la ligne précédente.}
\[
  \frac{4\pi}{5} < \frac{2\pi}{3} \qquad \frac{\pi}{3} < \qquad \frac{4\pi}{5} \to \frac{\pi}{6} \;? \qquad \frac{4\pi}{5} \leqslant \frac{5\pi}{6}
\]
\note{ces inégalités, éparses dans la marge des croquis, sont données telles quelles ; la première est fausse ($4\pi/5 > 2\pi/3$), la dernière a son $4$ récrit.}

$u, v', w'$ vecteurs unitaires
\[
  \begin{array}{l} u \perp v' \\ u \perp w' \end{array} \qquad v'w' = \cos\frac{2\pi}{5}
\]
\[
\begin{aligned}
  x(\alpha) &= u\cos\alpha + v'\sin\alpha \\
  y(\alpha) &= u\cos\alpha + w'\sin\alpha
\end{aligned}
\]
\[
  x(\alpha) - u = u(\cos\alpha - 1) + v'\sin\alpha
\]
\[
\begin{aligned}
  \|x(\alpha) - u\|^{2} = \|y(\alpha) - u\|^{2} &= (\cos\alpha - 1)^{2} + \sin^{2}\alpha = 2 - 2\cos\alpha \\
  &= 4\sin^{2}\frac{\alpha}{2} \qquad \Bigl(1 - \cos\alpha = 2\sin^{2}\frac{\alpha}{2}\Bigr)
\end{aligned}
\]
\[
  \|x(\alpha) - u\| = 2\sin\frac{\alpha}{2}
\]
\[
\begin{aligned}
  \|x(\alpha) - y(\alpha)\|^{2} &= \|v'\sin\alpha - w'\sin\alpha\|^{2} \\
  &= \sin^{2}\alpha \, \|v' - w'\|^{2} = \sin^{2}\alpha \Bigl(2 - 2\cos\frac{2\pi}{5}\Bigr) \\
  &\struck{= \sin^{2}\alpha + \sin^{2}\alpha - 2\cos\frac{2\pi}{5}\sin^{2}\alpha} \\
  &= 4\sin^{2}\alpha \sin^{2}\frac{\pi}{5} \struck{= \ill{}}
\end{aligned}
\]
\[
  \|x(\alpha) - y(\alpha)\| = 2\sin\alpha \sin\frac{\pi}{5} = 4\sin\frac{\alpha}{2}\cos\frac{\alpha}{2}\sin\frac{\pi}{5}
\]
\[
  \|x(\alpha) - u\| = \|y(\alpha) - u\| = \|x(\alpha) - y(\alpha)\|
\]
\[
  \Longleftrightarrow \quad 2\cos\frac{\alpha}{2}\sin\frac{\pi}{5} = 1 \quad \text{i.e.} \quad \boxed{\cos\frac{\alpha}{2} = \frac{1}{2\sin\pi/5}}
\]
\marginal{NB $2\sin\frac{\pi}{5}$ = côté du pentagone régulier \uncertain{unité}}

\page{18}
\note{feuillet de carnet à spirale ; le bloc de gauche est écrit en travers, le reste droit. Au centre, une projection de l'icosaèdre sur le plan orthogonal à un axe : un décagone dont les sommets portent alternativement $A_1, A'_3, A_0, A'_2, A_4, A'_1, A_2, A'_0, A_3, A'_4$ (plusieurs lettres récrites, $A$ par-dessus $\alpha$), le centre marqué $s$ et $s'$, les deux pentagones $A_i$ et $A'_i$ tracés, l'un en trait appuyé.}
$G$ groupe d'ordre 60, opérant \add{transitivement} sur $I$ ens. d'ordre 15, sur lequel il y a un automorphisme $\alpha$ d'ordre 3 commutant à $G$, d'où $I/\alpha = J$ ens. à 5 éléments, sur lequel $G$ opère transitivement. Le sous-groupe de stabilité d'un \add{$\mathrm{elt}$ de $J$} est un sous-groupe \struck{\uncertain{$\mathfrak{F}$}} d'ordre 12 de $G$, et on trouve $G \to \mathfrak{S}_5$ \add{ordre 120}. \ill{} \uncertain{si on sait} que $G$ opère fidèlement sur $J$, alors $G \to \mathfrak{S}_5$ est injectif, donc $G \simeq$ sous-groupe d'indice 2 de $\mathfrak{S}_5$, donc $G \simeq \mathfrak{A}_5$.
\note{« transitivement », « elt de $J$ » et « ordre 120 » sont interlinéaires.}
\[
\begin{aligned}
  (AB, A'B') &= P \\
  (EC', CE') &= Q = \alpha(P) \\
  (SD, S'D') &= R = \alpha(Q) = \alpha^{2}(P) \\
  (AB, A'B') &= \uncertain{\alpha(R)} = \alpha^{2}(Q) = \alpha^{3}(P)
\end{aligned}
\]
\note{les lettres $A, B, C, D, E, S$ ne sont pas portées sur la figure ; les couples sont apparemment des couples d'arêtes opposées. La dernière ligne est juste : $R = \alpha(Q)$ donne $\alpha(R) = \alpha^{2}(Q) = \alpha^{3}(P)$, et l'on revient à $P = (AB, A'B')$ parce que $\alpha$ est d'ordre 3.}

$\sigma \in G$, laissant stable $(P, Q, R)$, donc $\sigma \in N$ (normalisateur \uncertain{d'un} \ill{} de Sylow$_2$ \ill{}) ; l'intersection des conjugués de $N$ est un sous-groupe de $N$ invariant, donc soit $N$, soit $\uncertain{\mathfrak{F}_2}$, soit $0$. Ce n'est pas $N$, ni $\uncertain{\mathfrak{F}_2}$
\note{la phrase s'arrête au bord du feuillet.}

\[
  \alpha D \longrightarrow D \qquad \begin{array}{cc} 30 & 15 \\ \uncertain{\alpha^{!}}? & \alpha \end{array}
\]
$\alpha^{\uncertain{13}}$ est d'ordre 2 ou 1 [transform\ill{} \ill{} \ill{} \uncertain{orientation} \ill{} \ill{} \ill{}]
\[
  G/\mathfrak{F}_2 \longrightarrow G/\mathfrak{F}_4 \qquad 2, 4, 3, 1
\]
\note{au-dessus, quelques signes au crayon très pâle, illisibles.}
\[
  \frac{\pi}{\uncertain{3}} \leqslant \frac{2\pi}{5} \leqslant \frac{5\pi}{6}
  \qquad
  \cos\alpha = \cos^{2}\frac{\alpha}{2} - \sin^{2}\frac{\alpha}{2} = 2\cos^{2}\frac{\alpha}{2} - 1 = \frac{1}{2\sin^{2}\theta} - 1 = 1 - 2\sin^{2}\theta = \cos 2\theta
\]
\note{la chaîne pour $\cos\alpha$ est écrite en colonne, deux lignes intermédiaires se chevauchant ; on la donne telle qu'on la lit, sans en vérifier les égalités. Près d'elle, $\alpha'_{12} = \pi - \frac{\alpha}{2}$ (lecture douteuse), et un petit triangle rectangle dont l'hypoténuse porte $2\sin\frac{\alpha}{2}$ et un côté $4$, avec « $\ell^{2} = 4\sin^{2}\frac{\alpha}{2}$ ».}
\[
  u \to u\cos\alpha + v\sin\alpha \qquad v \to u(-\sin\alpha) + v\cos\alpha
\]

\page{19}
\note{feuillet de brouillon du même carnet, écrit dans deux sens (une partie à l'encre droite, la plus grande partie au crayon en travers), sans ordre suivi. On donne les formules lisibles par groupes ; les croquis (un triangle isocèle d'angle au sommet $\theta$ et de côtés $\rho$, avec $2\rho\sin\frac{\theta}{2}$ ; un cône de demi-angle $\alpha$ avec $\cos\alpha$ et un angle de $72^{\circ}$ ; trois vecteurs issus d'un point ; un décagone ; un cercle traversé d'un triangle rectangle) ne sont que signalés.}

\emph{À l'encre, en tête :}
\[
  \frac{1\,2\,3\,4\,\struck{0}}{2,4,1,3,0} \qquad \frac{}{3,1,4,2,0} \qquad \zeta \quad \zeta^{2}
\]
\[
  u, 2u, 3u, 4u, 0 \qquad \zeta^{i} \qquad e^{i\theta} \qquad u^{2} + w^{2} = 1 \qquad \theta^{2} + \mu^{2} = 1 \qquad x^{2} - y^{2}
\]
\note{les deux lignes de permutations $2,4,1,3,0$ et $3,1,4,2,0$ sont les images de $1,2,3,4,0$ par $i \mapsto 2i$ et $i \mapsto 3i$ modulo $5$ ; à côté, « $\mathfrak{S}_{i_0}$ » et une matrice commençant par $\cos\theta$, inachevée.}

\emph{Au crayon :}
\[
  g(\sigma . i) \qquad \rho(g)(\sigma . i) = \bigl(\rho(g)\ldots\bigr)\bigl(\rho(\ldots)\bigr) \qquad g(\sigma i) = g\sigma g^{-1}\, gi
\]
\[
  G/H = X \ni e \qquad y e = a . e = a \qquad y x = y g e = g \uncertain{u} e = g' a e = \ill{}
\]
\[
  b' = b h \qquad b^{-1} b' \in H \quad \Longrightarrow \quad (ba)^{-1} b' a \in H \quad \text{i.e.} \quad a^{-1} b^{-1} b' a \in H \quad \text{i.e.} \quad a^{-1} H a \subset H
\]
\[
  \frac{N(H')}{\bigcap \operatorname{Conj} H} \qquad
  N(H) \cap N(H') = N(\uncertain{H'}) \subset N(H) = \mathfrak{Z}_3 . \mathfrak{F}_2^{2}
  \qquad
  \begin{array}{ccc}
    N(H') & \subset & N(H) \\
    \cup & & \cup \\
    H' & \subset & H \;\text{---}\; \mathfrak{F}_2^{2} \\
    \cup & & \cup \\
    K \cap H' = K'_0 & \subset & K' = 0
  \end{array}
\]
\note{« $\mathfrak{Z}_3 . \mathfrak{F}_2^{2}$ » est récrit par-dessus d'autres signes ; la lecture est douteuse.}
\[
  \cos^{2}\alpha \cos\frac{2\pi}{5} = \cos\alpha \qquad 2\sin\frac{\alpha}{2}\sin\frac{2\pi}{5} = 1 \qquad \boxed{\cos\alpha = \frac{1}{2\sin\frac{2\pi}{\uncertain{10}}}}
\]
\[
  1 + x + x^{2} + x^{3} + x^{4} = 0 \qquad \Bigl(2\sin\frac{\alpha}{2}\Bigr)\Bigl(2\sin^{2}\frac{\pi}{\uncertain{4}}\Bigr) \qquad \alpha = 60\,?
\]
\[
   \cos\frac{\alpha}{2} = \frac{\sqrt{3}}{2} \qquad \sqrt{3}\sin\frac{\pi}{5} = 1 \qquad \frac{\pi}{6}
\]
\note{les relations entre $\alpha$ et $\pi/5$ de ce feuillet ne s'accordent pas toutes avec celle, encadrée, des pp. 4 et 17 ; ce sont des essais, transcrits sans correction.}

\end{document}
